% ==============================================================================
% sci_s1_overview.tex — SCI/SICE tutorial Session 1:
% overview and design philosophy of StampFly Ecosystem
% SCI/SICE チュートリアル セッション1: StampFly Ecosystem の全体像と設計思想
% ==============================================================================

% ------------------------------------------------------------------
\begin{frame}{地図: 今ここ}
  \begin{columns}[c]
    \begin{column}{0.5\textwidth}
      \centering
      \resizebox{\columnwidth}{!}{\scisessionmap{1}}
    \end{column}
    \begin{column}{0.47\textwidth}
      \small
      今日はこの循環を 5 つのセッションで 1 周する。まずは全体像と設計思想から
    \end{column}
  \end{columns}
\end{frame}

% ------------------------------------------------------------------
\begin{frame}{要点3つ}
  \begin{enumerate}
    \item マルチコプタは制御なしには1秒も飛べない。この性質が制御工学の教材として理想的である
    \item 必要なリソースが分散していることが普及の障壁になっている。StampFly Ecosystem は単一リポジトリへの集約でこれに応える
    \item ファームウェアは4階層のアクセス構造を持ち、初学者から研究者まで同じコードベースを共有しながら自分の関心に合わせて深さを選べる
  \end{enumerate}

  \vspace{0.3em}

  {\small 後半では、飛ばす前後で同じモデルを使って結果を確かめる仕組み（3 種類のシミュレーション）にも触れる}
\end{frame}

% ------------------------------------------------------------------
\begin{frame}{制御がなければ飛べない}
  \vspace{0.3em}

  \small マルチコプタは固定翼機のように滑空できず、ヘリコプタのようにオートローテーションで降りることもできない。制御系が止まれば即座に墜落する

  \vspace{0.5em}

  \begin{exampleblock}{教材としての魅力}
    \small この性質は安全上の課題であると同時に、学習者が「制御がなければ飛べない」という事実を身をもって体験できる特性でもある。フィードバック制御の役割を直感的に理解する入口になる
  \end{exampleblock}
\end{frame}

% ------------------------------------------------------------------
\begin{frame}{StampFly ハードウェア}
  \begin{columns}[T]
    \begin{column}{0.42\textwidth}
      \centering
      \includegraphics[width=\columnwidth]{Stampfly_onhand.jpg}
    \end{column}
    \begin{column}{0.56\textwidth}
      \begin{table}
        \centering\footnotesize
        \renewcommand{\arraystretch}{0.82}
        \setlength{\tabcolsep}{3pt}
        \begin{tabular}{lp{58mm}}
          \hline
          \textbf{項目} & \textbf{仕様} \\
          \hline
          MCU        & ESP32-S3（M5StampS3） \\
          IMU        & BMI270（6軸 400\,Hz） \\
          気圧       & BMP280 \\
          磁力計     & BMM150（3軸） \\
          測距       & VL53L3CX $\times$2（Bottom / Front） \\
          光学フロー & PMW3901 \\
          電流電圧   & INA3221 \\
          モータ     & コアレスモータ $\times$4 \\
          バッテリ   & 300\,mAh 1S LiPo \\
          質量       & 37\,g（バッテリ含む） \\
          コントローラ & M5Stack AtomS3 + Atom JoyStick（ESP-NOW） \\
          \hline
        \end{tabular}
      \end{table}
      \vspace{-2pt}

      {\scriptsize IMU（慣性計測装置）／ESP-NOW（ESP32 の無線プロトコル）。前方 ToF は未使用（高度推定は下方 ToF）}
    \end{column}
  \end{columns}
\end{frame}

% ------------------------------------------------------------------
\begin{frame}{4つの障壁}
  \small 手のひらサイズで安全性が高く安価な StampFly でも、要点②のリソース分散に加えて、教育現場での普及には次の4つの障壁がある

  \vspace{0.5em}

  \begin{itemize}\small\setlength{\itemsep}{4pt}
    \item \textbf{ブラックボックス問題}: 市販の教育用ドローンは飛ぶが、制御ループの中身に触れられない
    \item \textbf{理論と実装のギャップ}: PIDの数式は教科書で学べても、組込みのリアルタイム制御ループとして実装するには別のスキルが要る
    \item \textbf{体系的教材の不在}: センサ取得から状態推定、制御、実験、データ解析までを一貫して学べる教材が見当たらない
    \item \textbf{教育者の準備負担}: カリキュラム設計・教材作成・ツール整備を一人の教育者がすべて担うのは現実的でない
  \end{itemize}
\end{frame}

% ------------------------------------------------------------------
\begin{frame}{既存プラットフォームとの比較}
  \begin{table}
    \centering\footnotesize
    \begin{tabular}{lcccc}
      \hline
       & Crazyflie & Tello EDU\textsuperscript{*} & PX4 & StampFly Eco. \\
      \hline
      制御ループアクセス & \checkmark & --- & \checkmark & \checkmark \\
      教材・カリキュラム & $\triangle$ & $\triangle$ & $\triangle$ & \checkmark \\
      シミュレータ統合   & \checkmark & --- & \checkmark & \checkmark \\
      コスト             & 高 & 低 & 高 & 低 \\
      コード規模         & 中 & --- & 大 & 小 \\
      \hline
    \end{tabular}
  \end{table}

  \vspace{0.5em}

  {\small 制御ループへのフルアクセスと体系的な教材・ツールの統合提供を両立させ、ファームウェア全体を一人の学習者が読み通せる規模に抑えている点が異なる}

  {\scriptsize \textsuperscript{*}Tello EDU は2023年末に販売終了（評価は筆者の定性的見解）}
\end{frame}

% ------------------------------------------------------------------
\begin{frame}{エコシステムの地図}
  \centering
  \vspace{1mm}
  \resizebox{!}{0.78\textheight}{\input{../../_shared/tikz/ecosystem_layers}}
\end{frame}

% ------------------------------------------------------------------
\begin{frame}{5つのEcosystem設計の考え方}
  \begin{table}
    \centering\footnotesize
    \renewcommand{\arraystretch}{1.2}
    \setlength{\tabcolsep}{8pt}
    \begin{tabular}{lp{104mm}}
      \hline
      \textbf{考え方} & \textbf{内容} \\
      \hline
      一つのリポジトリ & ファーム・シミュレータ・解析ツール・教材を一か所にまとめる。通信仕様の定義は一つだけ置き、各実装はそこから作る \\
      sf CLI & ビルド・書込み・ログ取得・シミュレーションを一つのコマンド体系で操作する \\
      4 階層アクセス & 実習用 API からハードウェアの直接操作まで 4 つの入口を用意し、必要な深さから入る \\
      シミュレータ統合 & 軽量 3D・高精度物理・実ファーム実行の 3 種類のシミュレータで、同じ制御アルゴリズムを動かす \\
      オープンソース & 制御ループからドライバまで全部を公開し、ブラックボックスを作らない \\
      \hline
    \end{tabular}
  \end{table}
\end{frame}

% ------------------------------------------------------------------
\begin{frame}{1機から学べる8分野}
  {\small この 5 つの考え方の上で、1 機から次の 8 分野を扱える}

  \vspace{0.15em}

  \begin{table}
    \centering\footnotesize
    \renewcommand{\arraystretch}{1.0}
    \setlength{\tabcolsep}{10pt}
    \begin{tabular}{ll}
      \hline
      \textbf{分野} & \textbf{主な学習項目} \\
      \hline
      制御理論       & PID, カスケード制御, FF/FB, アンチワインドアップ \\
      飛行力学       & 6DoF運動方程式, ホバー条件, NED/Body座標変換 \\
      センサ工学     & IMU, 気圧, 光学フロー, Allan分散, ノイズ解析 \\
      状態推定       & 相補フィルタ, ESKF, センサフュージョン \\
      アクチュエータ & PWM制御, 推力モデル, プロペラ空力特性, ミキサ行列 \\
      システム同定   & 周波数応答, パラメータ推定, モデル検証 \\
      通信・IoT      & ESP-NOW, WiFiテレメトリ, プロトコル設計 \\
      ソフトウェア工学 & リアルタイムタスク, 組込み開発, OSS開発 \\
      \hline
    \end{tabular}
  \end{table}

  \vspace{0.15em}

  {\small これらは独立した知識ではなく、1台の機体の中で密接に関連し合っている。PIDゲインの設計には慣性モーメントやモータ応答特性の知識が要り、推定精度はセンサのノイズ特性に左右される}
\end{frame}

% ------------------------------------------------------------------
\begin{frame}{vehicle ファーム構造}
  \begin{table}
    \centering\footnotesize
    \begin{tabular}{ll}
      \hline
      コンポーネント数 & 30（\code{sf\_*} 命名, フラット構造） \\
      通信方式         & Pub-Sub トピック（直接呼び出し禁止） \\
      タスク数         & 16 \\
      制御周期         & 400\,Hz（IMU $\to$ 状態推定 $\to$ 制御） \\
      状態推定         & ESKF（誤差状態カルマンフィルタ） \\
      制御器           & カスケード PID \\
      飛行モード       & ACRO / STABILIZE / ALT\_HOLD / POS\_HOLD \\
      \hline
    \end{tabular}
  \end{table}

  {\footnotesize Pub-Sub（出版購読型のメッセージ配信）}

  \vspace{0.3em}

  \begin{exampleblock}{実機検証}
    \small POS\_HOLD（位置制御）は実機で検証済み。保持精度 $\pm$6--7\,cm
  \end{exampleblock}
\end{frame}

% ------------------------------------------------------------------
\begin{frame}{400\,Hz 制御ループ}
  \centering
  \vspace{-1mm}
  \resizebox{0.82\textwidth}{!}{\input{../../_shared/tikz/loop_400hz_timeline}}

  \vspace{0.3em}

  {\small IMU 割込みを起点に、状態推定 $\to$ 制御 $\to$ アクチュエーションが 2.5\,ms 周期で走る。全コンポーネントが同じ周期を共有する}
\end{frame}

% ------------------------------------------------------------------
\begin{frame}{4階層アクセス}
  \centering
  \vspace{-1mm}
  \resizebox{!}{0.82\textheight}{\input{../../_shared/tikz/sci_tier_access}}
\end{frame}

% ------------------------------------------------------------------
\begin{frame}{シミュレーション3層}
  {\small 設計の考え方「シミュレータ統合」を掘り下げる。目的の異なる 3 種類のシミュレーションを使い分ける}

  \vspace{0.15em}

  \begin{table}
    \centering\footnotesize
    \setlength{\tabcolsep}{5pt}
    \begin{tabular}{p{28mm}p{52mm}p{42mm}}
      \hline
      \textbf{層} & \textbf{内容} & \textbf{位置づけ} \\
      \hline
      層1 設計用線形モデル & 実飛行ログから同定した低次 $G(s)$ & 制御系設計・ループ整形の基準 \\
      層2 実ログ駆動再生   & 移植した制御則を実ログの外乱・指令で駆動 & パラメータ変更の A/B 判定の基準 \\
      層3 SILS             & MuJoCo + 実ファーム C++、Code/Param Identity & システム全体の検証ベンチ \\
      \hline
    \end{tabular}
  \end{table}

  \vspace{0.2em}

  {\footnotesize SILS（Software-in-the-Loop Simulation: 実ファームをホスト PC 上で MuJoCo 物理エンジンと組み合わせて動かすシミュレーション検証）}

  \vspace{0.1em}

  {\footnotesize 3 つとも同じ物理パラメータを使うので、飛ばす前の設計（線形モデル・ログ再生）と飛ばした後の答え合わせ（SILS）が同じモデルでつながる。合否の基準と実習は S5 で扱う}
\end{frame}

% ------------------------------------------------------------------
\begin{frame}{開発の3原則\later}
  \begin{table}
    \centering\footnotesize
    \begin{tabular}{lp{100mm}}
      \hline
      \textbf{原則} & \textbf{内容} \\
      \hline
      Code Identity  & \textbf{同じコードで動く。}SILS は実機と同じファームウェアのソースをそのまま組み込んで動かす。推定や制御の式だけでなく、制御ループ全体が実機と同一 \\
      Parameter Identity & \textbf{同じ設定値で動く。}ゲインなどの調整値は SILS と実機が同じ一覧から読む。片方で決めた値をそのまま他方へ持ち込める \\
      Model Fidelity & \textbf{正解は物理モデルから作る。}SILS の「真値」は物理モデルで作る。実機データが要るのは、初飛行のあとでモデルを現実に近づけるときだけ \\
      \hline
    \end{tabular}
  \end{table}
\end{frame}

% ------------------------------------------------------------------
% sf sils gui is listed in this session's review path, so show it once here
% (what the screen looks like). How to use it and the hands-on are in S5.
% 復習パスに sf sils gui が載るので、ここで 1 枚だけ画面を見せる。
% 使い方と実習は S5 で扱う。
\begin{frame}[fragile]{SILS を画面で触る: sf sils gui}
  \vspace{0.15em}

  \begin{columns}[T]
    \begin{column}{0.58\textwidth}
      \centering
      % Screenshot: pos_roll scenario just after a run (headless Chromium capture)
      % 画面写真: pos_roll シナリオ実行直後（ヘッドレス Chromium で撮影）
      \includegraphics[width=\columnwidth]{sils_gui.png}\\[0.2em]
      {\scriptsize \code{pos\_roll} シナリオ実行直後。左: イベント列、右上: 3D 再生、右下: グラフ（時刻カーソルが 3D と同期）、右上端: 合否 12/12 PASS}
    \end{column}
    \begin{column}{0.40\textwidth}
      \begin{itemize}\footnotesize\setlength{\itemsep}{1pt}
        \item ブラウザだけで完結: シナリオ作成（外乱・故障などのイベントを表で組む）$\to$ 実行 $\to$ グラフ $\to$ 合否判定
        \item PID ゲイン等 54 項目を画面で編集。再ビルド不要（Parameter Identity）
        \item 使い方と実習は S5。今日は画面を覚えるだけでよい
      \end{itemize}

      \vspace{0.2em}

      \begin{sflisting}[basicstyle=\ttfamily\scriptsize]{起動}
sf sils build
sf sils gui
      \end{sflisting}
    \end{column}
  \end{columns}
\end{frame}

% ------------------------------------------------------------------
\begin{frame}[fragile]{デモ①: シミュレータ操縦}
  \vspace{0.15em}

  {\small \textcolor{sfblue}{\textbf{見るだけ}} ｜ 一緒に打つ ｜ 帰宅後に再現}

  \vspace{0.2em}

  \begin{columns}[T]
    \begin{column}{0.54\textwidth}
      \begin{block}{何を見るか}
        \begin{itemize}\footnotesize\setlength{\itemsep}{0pt}
          \item \code{sf sim run vpython} を手元でも起動
          \item VPythonシミュレータが起動し、実機と同じ制御アルゴリズムが動く
          \item HIDジョイスティック（Atom JoyStick）でそのまま操縦できる
          \item 実機がなくても制御コードの挙動を3Dで確認できる
        \end{itemize}
      \end{block}
    \end{column}
    \begin{column}{0.44\textwidth}
      \centering
      % VPython simulator screenshot (VOXEL terrain, in flight); shared with the DXH deck
      % VPython シミュレータの画面（ボクセル地形・飛行中）。DXH 資料と共用
      \includegraphics[width=\columnwidth]{simulator_demo.png}\\[0.2em]
      {\scriptsize VPython シミュレータの画面（ボクセル地形を飛行中）}
    \end{column}
  \end{columns}
\end{frame}

% ------------------------------------------------------------------
\begin{frame}{デモ②: 実機 POS\_HOLD 飛行}
  \vspace{0.3em}

  {\small \textcolor{sfblue}{\textbf{見るだけ}} ｜ 一緒に打つ ｜ 帰宅後に再現}

  \vspace{0.5em}

  \begin{block}{何を見るか}
    \begin{itemize}\small
      \item ホバー中の位置保持精度（実測 $\pm$6--7\,cm）
      \item 手で軽く押すなどの外乱を与えたときの復帰動作
    \end{itemize}
  \end{block}
\end{frame}

% ------------------------------------------------------------------
\begin{frame}{デモ②の期待結果: 位置保持}
  \centering
  \vspace{0.3em}

  \includegraphics[width=\textwidth, height=0.5\textheight, keepaspectratio]{../fallback/S1_pos_hold_xy.png}

  {\scriptsize \code{pos\_roll.scn}（vehicle）: 離陸 $\to$ ロール外乱 $\to$ POS\_HOLD 係合 $\to$ 保持。係合後の水平ドリフト最大 0.39\,m}
  % Web版のみ: 静止グラフの代わりにPOS_HOLD飛行動画（MuJoCo 3D + 状態グラフ）を見せる
  % @web: video=../fallback/S1_pos_hold_flight.mp4 caption=POS_HOLD 飛行（MuJoCo 3D + 状態グラフ）
  % @web: layout=replace
\end{frame}

% ------------------------------------------------------------------
\begin{frame}{デモ②の期待結果: 高度と姿勢}
  \centering
  \vspace{0.3em}

  \includegraphics[width=\textwidth, height=0.6\textheight, keepaspectratio]{../fallback/S1_altitude_attitude.png}

  {\scriptsize 同じ飛行の高度と姿勢角の時系列。外乱直後の傾きが位置制御で戻る}
\end{frame}

% ------------------------------------------------------------------
\begin{frame}{復習パス\later}
  \begin{table}
    \centering\footnotesize
    \setlength{\tabcolsep}{4pt}
    \begin{tabular}{>{\raggedright\arraybackslash}p{35mm}>{\raggedright\arraybackslash}p{35mm}>{\raggedright\arraybackslash}p{65mm}}
      \hline
      \textbf{コマンド} & \textbf{対象} & \textbf{文書} \\
      \hline
      \code{sf sim run vpython} & シミュレータ操縦 & \code{docs/\allowbreak architecture/\allowbreak simulation-policy.md} \\
      \code{sf sils gui} & SILS シナリオ実行・可視化 & \code{simulator/\allowbreak sils/\allowbreak gui/\allowbreak README.md} \\
      --- & 4階層アクセス & \code{firmware/\allowbreak vehicle/\allowbreak docs/\allowbreak architecture.md \S2} \\
      \hline
    \end{tabular}
  \end{table}
\end{frame}

% ------------------------------------------------------------------
\begin{frame}{チェックポイント}
  \begin{alertblock}{確認事項}
    \begin{itemize}
      \item[$\square$] マルチコプタが制御教材として適している理由を説明できる
      \item[$\square$] StampFly Ecosystem が集約によってどの障壁に対応しているか説明できる
      \item[$\square$] 4階層アクセス（L0--L3）のうち自分の関心に合う層を1つ選べる
    \end{itemize}
  \end{alertblock}

  \vspace{1em}

  \begin{block}{次のセッション}
    S2: 開発環境のセットアップとセンサデータの取得（11:00 -- 12:00）--- 全体像を見た次は、実際に手を動かしてセンサの値を読む番
  \end{block}
\end{frame}
